针对空间站微重力环境下磁驱动柔性导管在磁场、流场与血管壁接触共同作用中的三维大变形动力学问题，本文构建从几何精确杆模型到结构化学习与 KAM 稳定性分析的统一框架。首先，以中心线位置和截面姿态为独立变量，建立保留伸长、剪切、弯曲、扭转及转动惯量的 Cosserat 杆模型，并通过能量与功率恒等式区分磁、流体、接触和边界作用中的保守项、耗散项与外部输入。其次，采用平移--转动联合 Galerkin 离散、约束流形局部约化和一致 Legendre 变换，获得有限维正则 Hamiltonian 系统及与 COMSOL 场数据一致的状态和端口接口。固定的 \(N_{\mathrm{diag}}=2\) 诊断基采用 \(q=\Pi_{\mathrm{conf}}(U)\)、\(p=B_N(q)^{T}P_{\mathrm{full}}\)、\(\dot q=M_N(q)^{-1}p\) 和 \(\dot p=\dot B_N(q)^{T}P_{\mathrm{full}}+B_N(q)^{T}\dot P_{\mathrm{full}}\) 的 paired-coordinate contract；该基用于检验数据映射的一致性，但不意味着 \(N_{\mathrm{diag}}=2\) 已被证明为最终充分的主模态阶数。COMSOL--Galerkin bridge 的三条候选轨迹均满足独立的 \(2\%\) force-v2 联合 L2 判据，但原 \(0.5\%\) 严格扰动能量判据总体不满足，两类判据不可相互替代。低维 P0 和三轨迹 P1/P1.5/P1.6 实验显示：P1 在导数监督的 pilot 上呈现条件性优势，P1.5 中灰箱 residual 相对 physics-only 明显改善但未形成稳定的一步场优势，P1.6 则表明精确满足平衡约束可能伴随预测误差显著增大。经跨轨迹比对修正、误差源隔离、余切投影合同识别、求解器输出语义比较和固定基初值一致性修正后，三层时间步比较得到 \(1.64186\)--\(2.00088\) 的观察收敛阶；\(1\,\mathrm{s}\) 幅值哨兵算例的 G2、G3 点态、G3 impulse 和扰动能量峰峰漂移分别为 \(0.00146226\)、\(0.000221162\)、\(0.0000563606\) 和 \(0.00307057<0.005\)。当前有限证据至多支持该单一诊断哨兵作为后续 A1 资格实验的候选输入；它不认证跨 mode、幅值或工况的数据生成资格，也不证明 BDF 为辛积分器或构成 P2、KAM 的数值认证。计划的 96 条轨迹数据集、完整规模 P2 训练、开放端口、高保真对照、数值不变环面与工程安全域仍需在后续研究中完成。
